Large coding agents are often evaluated as if their cost is determined only by model choice, yet long-running agent sessions repeatedly resend large stable prefixes: system instructions, tool schemas, repository policies, and continuity context. This report starts from a simple question: if model-side systems can make tokens cheaper through training and inference engineering, what can an individual agent builder do outside the model? We present \toolname, a Rust command-line tool and Codex skill that audits whether a local Codex configuration is prompt-cache friendly. The first implementation focuses on stable provider selection, Responses API usage, WebSocket session continuity, environment-key health, and stable model/reasoning settings for cache-aware router deployments. The central claim is deliberately narrow: the tool does not reduce context or change the agent's task; it reduces the amount of repeated context likely to be paid as uncached input. We propose an evaluation protocol that measures cached tokens, uncached paid input, estimated cost, latency, and task success side by side. The intended outcome is a practical testbed for verifying whether cache-aware harness engineering lowers repeated-input cost while preserving coding-agent quality.
